\documentclass[12pt,a4paper]{article}
\usepackage[utf8]{inputenc}
\usepackage[T1]{fontenc}
\usepackage[brazil]{babel}
\usepackage{lmodern}
\usepackage{microtype}
\usepackage[a4paper,margin=2.5cm]{geometry}
\usepackage{parskip}
\setlength{\parindent}{1.25cm}
\setlength{\parskip}{0.45em}
\begin{document}

\begin{center}
{\Large\bfseries Resenha crítica}\par
\vspace{0.8em}
\begin{minipage}{0.9\textwidth}
\small\textbf{Referência:} LI, Li; YU, Fengchao; ZHANG, Enting. \textit{A systematic review of learning task design for K-12 AI education: Trends, challenges, and opportunities}. Computers and Education: Artificial Intelligence, v. 6, 100217, 2024. DOI: 10.1016/j.caeai.2024.100217.
\end{minipage}
\end{center}



Segundo Li, Yu e Zhang (2024), ensinar inteligência artificial na Educação Básica é uma tarefa desafiadora e não se resume a colocar uma ferramenta digital nas mãos dos estudantes e supor que a aprendizagem acontecerá. O problema aparece quando se pensa no que será proposto durante a aula, no que o aluno precisa compreender e em quais sinais o professor poderá perceber que houve compreensão.

Os autores procuraram entender como a aprendizagem de inteligência artificial vem sendo trabalhada na educação K-12. Começaram com 1.036 trabalhos e, depois de selecionar aqueles que atendiam aos critérios da pesquisa, analisaram 47 estudos empíricos. O artigo não discute apenas uma tecnologia ou um programa específico. A atenção está voltada às tarefas usadas nas aulas, aos conceitos mobilizados pelos estudantes e aos dados que podem indicar se eles compreenderam aquilo que estudaram.

Esse recorte traz uma limitação que chama atenção para o contexto brasileiro. Como os autores trabalharam apenas com artigos publicados em inglês, muitas experiências realizadas em escolas brasileiras e divulgadas em português não aparecem na revisão. Ainda assim, o levantamento ajuda a enxergar problemas que também aparecem na nossa realidade escolar.

Grande parte dos estudos encontrados trata de \textit{machine learning} e utiliza oficinas, jogos, programação ou plataformas digitais. Isso mostra que a área tem apostado em atividades práticas. Entretanto, o próprio artigo indica que uma atividade no computador pode deixar a turma interessada sem garantir que os estudantes entenderam o conceito que está por trás dela. Mexer em uma tela, ajustar um controle ou conseguir uma resposta não é o mesmo que explicar por que aquele resultado apareceu.

Os autores também apontam que muitos estudos relatam resultados positivos, mas nem sempre mostram se os estudantes conseguiram usar o que aprenderam em outra situação. Para quem acompanha uma turma, essa diferença é importante. Um aluno pode acertar uma questão logo depois da explicação e ainda assim ter dificuldade quando encontra um problema novo, com outros dados ou outra pergunta. Por isso, questionários de opinião e testes rápidos ajudam, mas não resolvem sozinhos a avaliação da aprendizagem.

A leitura conversa diretamente com a pesquisa proposta para o Ensino Médio. A plataforma não foi pensada apenas para apresentar modelos de \textit{machine learning}. Em cada atividade, o estudante deve observar os dados disponíveis, compreender qual é o problema e registrar uma hipótese antes de alterar os controles. Depois, precisa explicar o que mudou. A regressão linear, a SVM, o k-NN, a árvore de decisão e o k-means permitem discutir situações diferentes, mas a meta não é decorar o nome dos modelos. O que interessa é acompanhar como os estudantes relacionam dados, regras, parâmetros e resultados.

Essa escolha também depende da presença do professor. O site pode mostrar uma reta, uma fronteira ou grupos de pontos, mas a discussão em sala é que ajuda a transformar essa visualização em explicação. Quando o estudante diz que uma previsão mudou, por exemplo, o professor pode perguntar qual dado foi usado, que controle foi alterado e por que o resultado mudou. São perguntas simples, mas elas impedem que a atividade fique restrita ao clique mecânico.

O artigo de Li, Yu e Zhang reforça a necessidade de testar propostas nas condições reais da escola pública. Tempo de aula, equipamentos disponíveis e ritmo da turma interferem no que pode ser feito. A sequência de oito encontros, o diagnóstico inicial, os registros durante as aulas e o desafio final foram pensados para acompanhar as explicações dos estudantes ao longo desse processo. O objetivo não é provar que a plataforma, sozinha, produz aprendizagem, mas compreender o que ela permite discutir quando é usada com uma sequência organizada e com intervenção docente.

O principal aprendizado trazido pelo artigo está no rompimento com a ideia de que a tecnologia resolve, por si só, os problemas de ensino. A ferramenta pode abrir possibilidades, mas é a atividade proposta, a pergunta feita pelo professor e a explicação construída pelo estudante que dão sentido ao uso dela na sala de aula.

\end{document}
